% BHU PYQ -- Manually transcribed & error-corrected
% Paper: BCH-312: Income Tax Law & Accounts
% Exam: B.Com. (Hons.) Semester V Examination, 2013-14

\documentclass[11pt,a4paper]{article}
\usepackage[a4paper,top=2.5cm,bottom=2.5cm,left=2.8cm,right=2.8cm,headheight=14pt]{geometry}
\usepackage{amsmath,amssymb}
\usepackage[T1]{fontenc}
\usepackage[utf8]{inputenc}
\usepackage{lmodern,microtype,fancyhdr,enumitem,hyperref,lastpage,parskip}

\hypersetup{
    colorlinks=true,
    urlcolor=black,
    linkcolor=black,
    pdftitle={BCH-312: Income Tax Law \& Accounts -- BHU B.Com. (Hons.) Sem V 2013-14}
}

\pagestyle{fancy}
\fancyhf{}
\renewcommand{\headrulewidth}{0.4pt}
\renewcommand{\footrulewidth}{0.4pt}
\fancyhead[L]{\small   \textit{Banaras Hindu University}}
\fancyhead[R]{\small   \textit{BCH-312: Income Tax Law \& Accounts}}
\fancyfoot[C]{\small Page  \thepage\ of \pageref{LastPage}}


\newlist{parts}{enumerate}{2}
\setlist[parts,1]{label=(\alph*),leftmargin=*,itemsep=5pt,topsep=3pt}
\setlist[parts,2]{label=(\roman*),leftmargin=*,itemsep=3pt,topsep=2pt}

\newcommand{\pts}[1]{\hfill   \textit{[#1]}}


\begin{document}


\begin{center}
  {\large\bfseries BANARAS HINDU UNIVERSITY}\\[2pt]
  {\normalsize B.Com. (Hons.) --- Semester V Examination, 2013-14}\\[6pt]
  \rule{0.8\linewidth}{0.5pt}\\[6pt]
  {\large\bfseries Paper No. BCH-312: Income Tax Law \& Accounts}\\[4pt]
  {\normalsize Time: 3 Hours\hfill Maximum Marks: 70}
\end{center}

\rule{\linewidth}{0.4pt}

\smallskip



\noindent   \textit{Instructions: Answer all questions. All questions carry equal marks (14 marks each).}

\bigskip




\noindent   \textbf{Question 1.} \\
Explain the terms ``Previous Year'' and ``Assessment Year''. What are the exceptions to the rule of previous year?\pts{14}

\centerline{   \textbf{OR}}

Explain the following terms under the Income Tax Act, 1961:\pts{3.5+3.5+3.5+3.5}

\begin{parts}
  \item Assessee
  \item Person
  \item Partly agricultural income
  \item Gross Total Income and Total Income
\end{parts}

\medskip\rule{\linewidth}{0.2pt}

\pagebreak




\noindent   \textbf{Question 2.} \\
Mr. X is employed in Bharatpur Company Ltd. The population of the city is 7 lakh. The following information is available in connection with his income for the year ending 31st March, 2013:\pts{14}

\begin{enumerate}[label=(\roman*),itemsep=2pt,topsep=2pt]
  \item Salary Rs. 10,000 p.m.
  \item City compensatory allowance @ Rs. 1,000 p.m.
  \item Bonus @ 8\% of basic pay.
  \item Employer contributes 15\% of his basic salary to recognized provident fund. X contributes an equal amount.
  \item A rent-free accommodation equipped with furniture has also been provided. The cost of furniture is Rs. 80,000 (depreciated value being Rs. 64,800).
  \item The employer company gifted him Rs. 15,000 on his marriage anniversary during the previous year.
  \item The employee has appointed a gardener and a watchman. They are paid by the company at Rs. 1,000 p.m. and Rs. 800 p.m. respectively.
  \item He is also getting an education allowance for his two children @ Rs. 450 p.m. for each.
  \item He paid Rs. 4,000 as professional tax for two years during the previous year.
\end{enumerate}
Ascertain the taxable salary of X for the assessment year 2013-14.

\centerline{   \textbf{OR}}

Mr. Smith is the owner of three houses, the municipal valuation of which is Rs. 20,000, Rs. 25,000 and Rs. 32,000 respectively. All these houses are used by Mr. Smith for his own residential purposes. Fair rent of all the three houses is generally 20\% higher than the municipal valuation. Annual rental value of all the three houses is Rs. 22,800, Rs. 27,000 and Rs. 36,000 respectively. He has furnished the following information:

\begin{enumerate}[label=(\roman*),itemsep=2pt,topsep=2pt]
  \item Municipal taxes are levied @ 10\% of municipal valuation.
  \item He has spent Rs. 4,000 on the repairs of the first house, Rs. 10,000 on interior decoration, painting, and repairs of the second house, and Rs. 100 on the repair of the third house.
  \item A watchman and sweeper are employed in each house. Salary of Rs. 8,000 is spent on each house.
  \item Fire insurance premium is 2\% of municipal valuation.
\end{enumerate}
Compute the taxable income of Mr. Smith under the head ``Income from House Property'' for the assessment year 2013-14.\pts{14}

\medskip\rule{\linewidth}{0.2pt}

\pagebreak




\noindent   \textbf{Question 3.} \\
Mr. Z commenced a business on 1st July, 2012 in Bengaluru. The final accounts of the business were prepared on 31st March, 2013. His Profit \& Loss Account was as under:\pts{14}

\begin{center}
\small
\renewcommand{\arraystretch}{1.2}

\begin{tabular}{|l|r|l|r|}
\hline
   \textbf{Particulars} &   \textbf{Amount (Rs.)} &   \textbf{Particulars} &   \textbf{Amount (Rs.)} \\ \hline
Advertisement & 10,000 & Gross Profit & 1,00,000 \\
Preliminary Expenses & 2,500 & Interest on securities & 6,400 \\
Expenses on Scientific Research & 10,000 & Income from house property & 6,000 \\
Commission & 12,500 & & \\
Provision for Depreciation & 6,000 & & \\
Provision for Bad Debts & 2,000 & & \\
Embezzlement by employee & 2,500 & & \\
Salaries & 10,000 & & \\
Office Expenses & 5,000 & & \\
Discount to dealers & 2,500 & & \\
Investment allowance reserve & 7,500 & & \\
Sundry Expenses & 7,500 & & \\
Net Profit & 34,400 & & \\ \hline
   \textbf{Total} &   \textbf{1,12,400} &   \textbf{Total} &   \textbf{1,12,400} \\ \hline
\end{tabular}
\end{center}




\noindent   \textbf{Additional Information:}

\begin{enumerate}[label=(\roman*),itemsep=2pt,topsep=2pt]
  \item Total cost of project is Rs. 6,00,000.
  \item Expenses on scientific research include Rs. 5,000 in respect of building constructed for laboratory.
  \item Sundry Expenses include Rs. 5,000 given to a factory inspector as a bribe.
  \item An income of Rs. 2,000 has been earned but has not been credited to the Profit \& Loss Account.
  \item Sundry Expenses include Rs. 1,000, being the cost of a neon sign board.
  \item Advertisement expenses include Rs. 1,000 in respect of an advertisement for the loss of papers of Mr. Z's building.
\end{enumerate}
Compute the income of Mr. Z from Business and Profession for the assessment year 2013-14.

\centerline{   \textbf{OR}}

Explain the provisions of the Income Tax Act with regard to the following:\pts{5+5+4}

\begin{parts}
  \item Block of assets
  \item Normal depreciation and additional depreciation
  \item Unabsorbed depreciation
\end{parts}

\medskip\rule{\linewidth}{0.2pt}

\pagebreak




\noindent   \textbf{Question 4.} \\
What is `Cost of acquisition' of capital assets? How is the cost of acquisition determined under different conditions?\pts{14}

\centerline{   \textbf{OR}}

Explain the provisions of Section 80C and Section 80G of the Income Tax Act, 1961.\pts{14}

\medskip\rule{\linewidth}{0.2pt}




\noindent   \textbf{Question 5.} \\
Discuss any five instances when an assessee has a statutory obligation to deduct tax at source (TDS) and the consequences of non-compliance.\pts{14}

\centerline{   \textbf{OR}}

What are the various tax authorities envisaged in the Income Tax Act, 1961? Briefly outline the powers of tax authorities.\pts{14}

\vfill
\rule{\linewidth}{0.4pt}

\begin{center}\small
\href{https://bhu-exam.vercel.app/}{Click here to visit our website}\\[4pt]
\href{https://me-aryan.vercel.app/}{Click here to visit our contributor}\\[4pt]
\href{https://quashan-me.vercel.app/}{Click here to visit our contributor}
\end{center}

\end{document}
